% Part: lambda-calculus
% Chapter: introduction
% Section: term-revisited

\documentclass[../../../include/open-logic-section]{subfiles}

\begin{document}

\olfileid{lam}{syn}{tr}
\olsection{$\alpha$-সমতুল্যতা-শ্রেণি হিসেবে পদ}

এখন থেকে আমরা পদগুলিকে $\alpha$-সমতুল্যতা পর্যন্ত অভিন্ন বলে বিবেচনা
করব। অর্থাৎ, আমরা যখন কোনো পদ লিখব, তখন বোঝাব পদটি যে
$\alpha$-সমতুল্যতা-শ্রেণিতে আছে, সেই শ্রেণিকে। যেমন,
$\lambd[a][\lambd[b][a c]]$ লিখে আমরা তার সঙ্গে $\alpha$-সমতুল্য সব
পদের সেট বোঝাব; এর মধ্যে আছে $\lambd[a][\lambd[b][a c]]$,
$\lambd[b][\lambd[a][b c]]$ ইত্যাদি।

আরও, আগের অনুচ্ছেদগুলিতে $N, Q$-এর মতো অক্ষর দিয়ে কোনো পদ বোঝানো
হত; এখন থেকে সেগুলি দিয়ে একটি শ্রেণি বোঝাব, এবং পরবর্তী আলোচনায়
পদের বদলে এই শ্রেণিগুলিই হবে আমাদের অধ্যয়নের বিষয়। $x, y$-এর মতো
অক্ষরগুলি আগের মতোই চলরাশি বোঝাবে।

শ্রেণি~$M$-এর একটি নির্বিচার !!{element} বোঝাতে আমরা $\rep{M}$ সংকেতও
গ্রহণ করি; একাধিক প্রতিনিধি দরকার হলে লিখি $\rep{M}[0]$,
$\rep{M}[1]$ ইত্যাদি।
% BN-SRC-285 TARGET-CORRECTION-BEGIN
\par\smallskip\noindent\textnormal{\small\textbf{উৎস-সংশোধন (BN-SRC-285):} হিমায়িত ইংরেজি বাক্যে দুই প্রতিনিধি-সংকেতের সঙ্গে “etc.”-কেও একই গণিত-পরিসরের ভিতরে রাখা হয়েছে, ফলে গদ্য-সংক্ষিপ্তিটি ভুলভাবে গাণিতিক পাঠ্য হয়। বাংলা বাক্যে প্রতিনিধি-সংকেত দুটিকে আলাদা গণিত-পরিসরে রেখে “ইত্যাদি” গদ্যে রাখা হয়েছে।} % BN-SRC-285 TARGET-CORRECTION-END

বর্ণনা সহজ রাখতে পদের সংকেতগুলিই আবার ব্যবহার করি। শ্রেণিগুলির উপর
নিচের সংজ্ঞা নিই:
\begin{defn}
  \begin{enumerate}
  \item $\lambd[x][N]$ হলো $\lambd[x][\rep{N}]$-কে ধারণকারী শ্রেণি।
  \item $PQ$ হলো $\rep{P}\rep{Q}$-কে ধারণকারী শ্রেণি।
  \end{enumerate}
\end{defn}

এগুলি যে সুসংজ্ঞায়িত, তা দেখা কঠিন নয়; কারণ $\alpha$-রূপান্তর
সামঞ্জস্যশীল।

\begin{defn} \ollabel{def:fv}
  কোনো $\alpha$-সমতুল্যতা-শ্রেণি~$M$-এর \emph{মুক্ত চলরাশির সেট},
  অর্থাৎ $\FV{M}$, হলো $\FV{\rep{M}}$।
\end{defn}

এটি সুসংজ্ঞায়িত; কারণ \olref[alp]{thm:fv}-এ দেখানো হয়েছে যে
$\FV{\rep{M}[0]} = \FV{\rep{M}[1]}$।

শ্রেণির মধ্যে প্রতিস্থাপনের জন্যও একই সংকেত আবার ব্যবহার করি:
\begin{defn} \ollabel{def:sub}
  $M$-এ $y$-এর জায়গায় $R$-এর \emph{প্রতিস্থাপন}, অর্থাৎ
  $\Subst{M}{R}{y}$, হলো $\Subst{\rep{M}}{\rep{R}}{y}$-কে ধারণকারী
  শ্রেণি; এখানে এমন যেকোনো $\rep{M}$ ও~$\rep{R}$ নেওয়া যায় যাতে
  প্রতিস্থাপনটি সংজ্ঞায়িত।
% BN-SRC-287 TARGET-CORRECTION-BEGIN
\par\smallskip\noindent\textnormal{\small\textbf{উৎস-সংশোধন (BN-SRC-287):} সংজ্ঞাটির বাম পাশ একটি $\alpha$-সমতুল্যতা-শ্রেণি, কিন্তু হিমায়িত উৎস সেটিকে সরাসরি প্রতিনিধি-স্তরের পদ $\Subst{\rep{M}}{\rep{R}}{y}$-এর সমান বলে। পরের বাক্যের \olref[alp]{cor:sub} শুধু প্রমাণ করে যে ভিন্ন বৈধ প্রতিনিধি নিলে ফলগুলি $\alpha$-সমতুল্য। তাই বাংলা সংজ্ঞায় বাম পাশকে ওই প্রতিনিধি-স্তরের ফলটি যে শ্রেণিতে আছে, সেই শ্রেণি বলা হয়েছে।} % BN-SRC-287 TARGET-CORRECTION-END
\end{defn}

\olref[alp]{cor:sub} অনুযায়ী এটিও সুসংজ্ঞায়িত।

এই সংজ্ঞা আমাদের যুক্তিকে কতটা সহজ করে, লক্ষ করো। যেমন:
\begin{align}
  \Subst{\lambd[x][x]}{y}{x} & \ollabel{eq:1}\\
  &= \Subst{\lambd[z][z]}{y}{x} \ollabel{eq:2}\\
  &= \lambd[z][\Subst{z}{y}{x}] \\
  &= \lambd[z][z]
\end{align}
% BN-SRC-288 TARGET-CORRECTION-BEGIN
\par\smallskip\noindent\textnormal{\small\textbf{উৎস-সংশোধন (BN-SRC-288):} হিমায়িত সারিবদ্ধ গণনার প্রথম সারিতে $\Subst{\lambd[x][x]}{y}{x}$-এর পরে একটি সমতা-চিহ্ন আছে, কিন্তু তার ডান পাশে কোনো অভিব্যক্তি নেই; পরের সারি আবার সমতা-চিহ্ন দিয়েই শুরু হয়। বাংলা গণনায় প্রথম সারিটি শুধু \olref{eq:1}-এ চিহ্নিত করা হয়েছে এবং পরের সারির বৈধ সমতা অপরিবর্তিত রাখা হয়েছে।} % BN-SRC-288 TARGET-CORRECTION-END

এগুলিকে এখনও পদের উপর প্রতিস্থাপন মনে করলে \olref{eq:1}
অসংজ্ঞায়িত। কিন্তু আগেই বলা হয়েছে, এখন আমরা শ্রেণির উপর প্রতিস্থাপন
বিবেচনা করছি। সেই কারণেই \olref{eq:2} সম্ভব: একই শ্রেণিভুক্ত বলে
$\lambd[x][x]$-এর বদলে $\lambd[z][z]$ নেওয়া যায়।

একই কারণে এখন থেকে ধরে নেব, আমাদের বেছে নেওয়া প্রতিনিধিগুলি
প্রতিস্থাপনের জন্য প্রয়োজনীয় শর্তগুলি সর্বদা পূরণ করে। যেমন,
$\Subst{\lambd[x][N]}{R}{y}$ দেখলে ধরে নেব, প্রতিনিধি
$\lambd[x][N]$ এমনভাবে বেছে নেওয়া হয়েছে যাতে $x \neq y$ এবং
$x \notin \FV{R}$।
% BN-SRC-286 TARGET-CORRECTION-BEGIN
\par\smallskip\noindent\textnormal{\small\textbf{উৎস-সংশোধন (BN-SRC-286):} এই অনুচ্ছেদে হিমায়িত উৎস মুক্ত-চলরাশির জন্য সংজ্ঞায়িত $\FV{\cdot}$ সংকেতের বদলে পাঁচটি কাঁচা $FV(\cdot)$ অভিব্যক্তি ব্যবহার করে। বাংলা সংজ্ঞা, সুসংজ্ঞায়িততার সমীকরণ এবং প্রতিনিধির সতেজতা-শর্তে সব কটিকে অধ্যায়ের সংজ্ঞায়িত $\FV{\cdot}$ সংকেতে স্বাভাবিক করা হয়েছে।} % BN-SRC-286 TARGET-CORRECTION-END

$\lambd[x][x]$-কে “শ্রেণি” বলা কিছুটা অস্বাভাবিক শোনায়। তাই এখন
থেকে এগুলিকে $\Lambda$-পদ (অথবা এই অংশের বাকি জায়গায় শুধু “পদ”)
বলব; এতে আমাদের পরিচিত $\lambda$-পদগুলির থেকে এগুলিকে আলাদা করা যাবে।

\begin{editorial}
  এখনই পদকে বিদায় বলা যাচ্ছে না: $\Lambda$-পদের গোটা সংজ্ঞাই
  $\lambda$-পদের উপর ভিত্তি করে, এবং $\Lambda$-পদের উপর অপেক্ষক
  সংজ্ঞায়িত করার কোনো পদ্ধতি আমরা এখনও দিইনি। ফলে এমন সব অপেক্ষককে
  আগে $\lambda$-পদের উপর সংজ্ঞায়িত করতে হয়, তারপর প্রতিস্থাপনের মতো
  করে $\Lambda$-পদে “প্রক্ষেপণ” করতে হয়। তবে পাঠক স্বজ্ঞাতভাবে বুঝতে
  পারবেন, $\Lambda$-পদের উপর অপেক্ষক কীভাবে সংজ্ঞায়িত করা যায়।
\end{editorial}
\end{document}
